\chapter{Análise de Dados}

\section{Fundamentos de análise de dados}

Análise de Dados transforma observações em informação para descrever fenômenos, testar hipóteses, estimar relações, detectar padrões e apoiar decisões.

Uma análise pode ser:
\begin{itemize}
\item \textbf{descritiva}: o que ocorreu?;
\item \textbf{diagnóstica}: quais fatores estão associados ao ocorrido?;
\item \textbf{preditiva}: o que tende a ocorrer?;
\item \textbf{prescritiva}: qual ação otimiza determinado objetivo sob restrições?
\end{itemize}

Predição e causalidade são objetivos diferentes. Um modelo pode prever muito bem sem identificar a causa do fenômeno.

\section{População, amostra e unidade de análise}

\termo{População} é conjunto de interesse. \termo{Amostra} é subconjunto observado. \termo{Unidade de análise} é entidade sobre a qual cada observação se refere.

Antes de calcular qualquer estatística, é necessário responder: \emph{uma linha representa o quê?}

Em base de pagamentos, uma linha pode ser pagamento; em base de contratos, uma linha pode ser contrato. Juntar as duas sem respeitar granularidade pode duplicar valores.

\section{Tipos de variáveis}

\subsection{Qualitativas}

Nominais representam categorias sem ordem intrínseca, como órgão. Ordinais possuem ordem, como classificação baixo/médio/alto.

\subsection{Quantitativas}

Discretas assumem valores contáveis, como número de processos. Contínuas podem assumir valores em intervalo, como tempo de atendimento.

A escala da variável determina estatísticas e gráficos adequados.

\section{Medidas de tendência central}

\subsection{Média}

Para $n$ observações:
\[
\bar{x}=\frac{1}{n}\sum_{i=1}^{n}x_i.
\]

A média utiliza todos os valores e é sensível a extremos.

\subsection{Mediana}

É o valor central após ordenar dados. Para quantidade par, utiliza-se a média dos dois valores centrais segundo definição usual.

É robusta a outliers.

\subsection{Moda}

É o valor mais frequente. Pode existir mais de uma moda ou nenhuma única moda.

\begin{exemplo}[média e mediana]
Tempos em dias: $2,3,3,4,28$.

Média:
\[
\frac{2+3+3+4+28}{5}=8.
\]
Mediana = 3.

O valor 28 desloca fortemente a média. Para descrever o tempo ``típico'', a mediana pode ser mais representativa, sem eliminar a importância do caso extremo.
\end{exemplo}

\section{Quartis, percentis e boxplot}

Percentil $p$ divide distribuição de modo que aproximadamente $p\%$ das observações estejam abaixo ou iguais ao valor, conforme método de interpolação adotado.

Quartis:
\begin{itemize}
\item $Q_1$: 25º percentil;
\item $Q_2$: mediana;
\item $Q_3$: 75º percentil.
\end{itemize}

Intervalo interquartil:
\[
IQR=Q_3-Q_1.
\]

Regra comum de boxplot marca como potenciais outliers valores abaixo de
\[
Q_1-1{,}5IQR
\]
ou acima de
\[
Q_3+1{,}5IQR.
\]

Isso é heurística, não prova de erro ou fraude.

\section{Dispersão}

\subsection{Amplitude}

\[
R=x_{max}-x_{min}.
\]
É simples e extremamente sensível a extremos.

\subsection{Variância}

Variância populacional:
\[
\sigma^2=\frac{1}{N}\sum_{i=1}^{N}(x_i-\mu)^2.
\]

Variância amostral:
\[
s^2=\frac{1}{n-1}\sum_{i=1}^{n}(x_i-\bar{x})^2.
\]

O denominador $n-1$ corrige viés da estimativa da variância populacional sob condições usuais.

\subsection{Desvio padrão}

\[
s=\sqrt{s^2}.
\]
Possui a mesma unidade dos dados.

\subsection{Coeficiente de variação}

\[
CV=\frac{s}{\bar{x}}\times100\%.
\]
Permite comparar dispersão relativa entre conjuntos com escalas semelhantes e médias positivas adequadas. Quando média é próxima de zero, CV pode ser instável ou sem sentido.

\section{Padronização}

Z-score:
\[
z=\frac{x-\mu}{\sigma}.
\]

Valor $z=2$ indica observação dois desvios-padrão acima da média.

Padronização não transforma qualquer distribuição em normal; apenas centraliza e escala.

\section{Probabilidade}

Para eventos $A$ e $B$:
\[
P(A^c)=1-P(A).
\]

Regra da união:
\[
P(A\cup B)=P(A)+P(B)-P(A\cap B).
\]

Se eventos são mutuamente exclusivos, $P(A\cap B)=0$.

\section{Probabilidade condicional}

\[
P(A\mid B)=\frac{P(A\cap B)}{P(B)},\quad P(B)>0.
\]

A probabilidade de A dado B pode diferir muito da probabilidade marginal de A.

\subsection{Independência}

A e B são independentes se:
\[
P(A\cap B)=P(A)P(B).
\]
Equivalentemente, quando probabilidades são positivas:
\[
P(A\mid B)=P(A).
\]

Mutuamente exclusivo não significa independente. Se dois eventos não podem ocorrer juntos e ambos têm probabilidade positiva, saber que um ocorreu elimina o outro, logo há dependência.

\section{Teorema de Bayes}

\[
P(A\mid B)=\frac{P(B\mid A)P(A)}{P(B)}.
\]

\begin{exemplo}[base rate]
Uma irregularidade ocorre em 1\% dos contratos. Um teste tem sensibilidade 90\% e taxa de falso positivo 5\%.

Em 10.000 contratos:
\begin{itemize}
\item 100 irregulares $\Rightarrow$ 90 positivos verdadeiros;
\item 9.900 regulares $\Rightarrow$ 495 falsos positivos.
\end{itemize}

Entre 585 alertas, apenas 90 são positivos reais:
\[
PPV=\frac{90}{585}\approx15{,}4\%.
\]
Alta sensibilidade não implica que a maioria dos alertas seja verdadeira quando a prevalência é baixa.
\end{exemplo}

\section{Variáveis aleatórias e valor esperado}

Uma variável aleatória associa resultados a números.

Valor esperado discreto:
\[
E[X]=\sum_x xP(X=x).
\]

Variância:
\[
Var(X)=E[(X-E[X])^2]=E[X^2]-E[X]^2.
\]

\section{Distribuição Bernoulli}

Modela um ensaio com dois resultados, sucesso com probabilidade $p$ e fracasso com $1-p$.

\[
E[X]=p,\qquad Var(X)=p(1-p).
\]

\section{Distribuição binomial}

Se $X$ conta sucessos em $n$ ensaios Bernoulli independentes com mesma probabilidade $p$:
\[
P(X=k)=\binom{n}{k}p^k(1-p)^{n-k}.
\]

\[
E[X]=np,\qquad Var(X)=np(1-p).
\]

\begin{exemplo}
Se cada requisição falha independentemente com $p=0{,}02$, em 10 requisições a probabilidade de exatamente uma falha é:
\[
\binom{10}{1}(0{,}02)(0{,}98)^9\approx0{,}1667.
\]
\end{exemplo}

\section{Distribuição de Poisson}

Modela contagem de eventos em intervalo sob hipóteses específicas, com taxa média $\lambda$:
\[
P(X=k)=e^{-\lambda}\frac{\lambda^k}{k!}.
\]

\[
E[X]=Var(X)=\lambda.
\]

É útil como modelo inicial para chegadas, falhas ou eventos raros, mas dados reais podem apresentar sobredispersão e dependência.

\section{Distribuição normal}

A normal é simétrica e determinada por média $\mu$ e desvio padrão $\sigma$.

A padronização converte:
\[
Z=\frac{X-\mu}{\sigma}.
\]

A regra aproximada 68--95--99,7 indica proporção em 1, 2 e 3 desvios-padrão para distribuição normal.

Não se deve aplicar z-score como detector universal de outlier quando distribuição é fortemente assimétrica.

\section{Amostragem}

\subsection{Aleatória simples}

Cada unidade possui probabilidade conhecida e apropriada de seleção, sob desenho simples.

\subsection{Estratificada}

População é dividida em estratos e amostras são selecionadas em cada grupo. Pode aumentar precisão e garantir representação de subgrupos.

\subsection{Por conglomerados}

Selecionam-se grupos naturais e observam-se unidades dentro deles. É operacionalmente eficiente, mas correlação intragrupo pode aumentar variância.

\subsection{Sistemática}

Seleciona-se ponto inicial e elementos em intervalos regulares. Periodicidade na lista pode introduzir viés.

\section{Viés de seleção}

Amostra grande não corrige desenho enviesado.

Exemplo: analisar apenas contratos já auditados para estimar taxa de irregularidade de todos os contratos pode superestimar prevalência se auditorias foram direcionadas aos casos suspeitos.

\section{Teorema Central do Limite}

Sob condições adequadas, distribuição da média amostral aproxima-se de normal conforme tamanho da amostra cresce, independentemente da distribuição original.

Erro padrão da média:
\[
SE(\bar X)=\frac{\sigma}{\sqrt{n}}
\]
ou estimado por $s/\sqrt{n}$.

Dobrar amostra não reduz erro padrão pela metade. Para reduzir SE pela metade, é necessário aproximadamente quadruplicar $n$.

\section{Intervalo de confiança}

Com variância conhecida ou aproximação normal apropriada:
\[
\bar{x}\pm z_{\alpha/2}\frac{\sigma}{\sqrt{n}}.
\]

Para 95\%, $z\approx1{,}96$.

\begin{atencao}
Depois que um intervalo frequentista específico foi calculado, não é correto dizer que há 95\% de probabilidade de o parâmetro fixo estar naquele intervalo. O procedimento, repetido muitas vezes, cobre o parâmetro em aproximadamente 95\% das amostras sob as hipóteses.
\end{atencao}

\section{Teste de hipóteses}

Define-se hipótese nula $H_0$ e alternativa $H_1$.

\subsection{Erro Tipo I}

Rejeitar $H_0$ quando é verdadeira. Probabilidade controlada pelo nível de significância $\alpha$.

\subsection{Erro Tipo II}

Não rejeitar $H_0$ quando alternativa relevante é verdadeira. Probabilidade $\beta$.

Poder:
\[
1-\beta.
\]

Aumentar amostra tende a aumentar poder para detectar determinado efeito.

\section{p-valor}

p-valor é probabilidade, assumindo $H_0$ e o modelo, de obter estatística tão extrema ou mais extrema que a observada.

Não é:
\begin{itemize}
\item probabilidade de $H_0$ ser verdadeira;
\item tamanho do efeito;
\item probabilidade de resultado ter ocorrido por acaso em sentido genérico.
\end{itemize}

p pequeno pode acompanhar efeito irrelevante em amostra gigantesca. Significância prática precisa ser avaliada separadamente.

\section{Correlação}

Coeficiente de Pearson mede associação linear:
\[
r=\frac{Cov(X,Y)}{s_Xs_Y}.
\]

$r$ varia entre -1 e 1.

Correlação zero não implica independência: relação pode ser não linear.

\subsection{Spearman}

Spearman calcula correlação sobre ranks e captura associação monotônica, sendo menos sensível a outliers e escala linear.

\section{Correlação não implica causalidade}

Uma associação pode surgir por:
\begin{itemize}
\item causalidade X $\rightarrow$ Y;
\item causalidade Y $\rightarrow$ X;
\item variável de confusão Z;
\item seleção;
\item coincidência;
\item viés de medição.
\end{itemize}

Exemplo: órgãos maiores podem ter mais contratos e mais incidentes. Correlação entre número de contratos e incidentes não prova que contratos causam incidentes; tamanho do órgão pode explicar ambos.

\section{Regressão linear simples}

Modelo:
\[
Y=\beta_0+\beta_1X+\varepsilon.
\]

OLS escolhe coeficientes que minimizam soma dos quadrados dos resíduos:
\[
\sum_i(y_i-\hat{y}_i)^2.
\]

$\beta_1$ representa mudança média esperada em Y associada a uma unidade de X, sob o modelo. Em regressão observacional, isso não deve ser interpretado automaticamente como efeito causal.

\section{Regressão múltipla}

\[
Y=\beta_0+\beta_1X_1+\cdots+\beta_pX_p+\varepsilon.
\]

Cada coeficiente representa associação de sua variável mantendo as demais constantes no modelo.

\subsection{Multicolinearidade}

Preditores fortemente correlacionados aumentam incerteza dos coeficientes e dificultam interpretação. VIF é diagnóstico comum.

\section{Hipóteses da regressão}

Para inferência OLS clássica, são importantes linearidade adequada, independência/estrutura de erros, homocedasticidade e condições sobre distribuição para certos testes.

Violação não torna automaticamente previsão impossível, mas afeta interpretação e inferência.

\section{R² e R² ajustado}

\[
R^2=1-\frac{SS_{res}}{SS_{tot}}.
\]

R² mede proporção de variação explicada em relação à média dentro da amostra/modelo. Pode aumentar ao adicionar variáveis, mesmo irrelevantes.

R² ajustado penaliza quantidade de preditores.

R² alto não prova causalidade nem bom desempenho fora da amostra.

\section{Regressão logística}

Para classificação binária:
\[
P(Y=1\mid X)=\sigma(z)=\frac{1}{1+e^{-z}},
\]
com
\[
z=\beta_0+\beta_1X_1+\cdots+\beta_pX_p.
\]

Log-odds são lineares:
\[
\log\frac{p}{1-p}=\beta_0+\beta_1X_1+\cdots.
\]

$e^{\beta_j}$ é razão de odds associada a incremento unitário, mantendo outros preditores constantes.

\section{Matriz de confusão}

\begin{center}
\begin{tabular}{c|cc}
 & Predito positivo & Predito negativo\\\hline
Real positivo & TP & FN\\
Real negativo & FP & TN
\end{tabular}
\end{center}

Métricas:
\[
Accuracy=\frac{TP+TN}{TP+TN+FP+FN},
\]
\[
Precision=\frac{TP}{TP+FP},
\]
\[
Recall=\frac{TP}{TP+FN},
\]
\[
Specificity=\frac{TN}{TN+FP}.
\]

\section{F1-score}

\[
F_1=2\frac{Precision\cdot Recall}{Precision+Recall}.
\]

É média harmônica e penaliza forte desequilíbrio entre precisão e recall.

Não considera TN diretamente, portanto não é métrica universal.

\section{Classes desbalanceadas}

Se apenas 1\% dos casos é positivo, classificador que prevê sempre negativo tem 99\% de acurácia e recall zero.

Métricas devem refletir objetivo e custos.

\begin{exemplo}[capacidade operacional]
Existem 10.000 contratos e unidade só pode investigar 100 por mês. Um modelo com recall alto, mas que gera 4.000 alertas, pode ser inviável. Métricas como precision@100 ou recall dentro do top-k podem estar mais alinhadas à decisão.
\end{exemplo}

\section{Threshold}

Modelos probabilísticos produzem score; threshold converte score em decisão.

Reduzir threshold aumenta, em geral, recall e também falsos positivos. Ele não é propriedade fixa do modelo e deve ser governado segundo custo dos erros.

\section{ROC e AUC}

ROC relaciona True Positive Rate a False Positive Rate ao variar threshold.

AUC pode ser interpretada como probabilidade de um positivo aleatório receber score maior que um negativo aleatório, sob condições usuais e ausência de empates tratados.

AUC mede ranking, não calibração.

\section{Precision-Recall Curve}

Em classes raras, PR curve destaca relação entre precisão e recall dos positivos e pode ser mais informativa que ROC.

Baseline de precisão depende da prevalência.

\section{Calibração}

Modelo calibrado com score 0,8 deve observar aproximadamente 80\% de positivos entre casos semelhantes com score próximo a 0,8.

Brier score:
\[
\frac{1}{n}\sum_i(p_i-y_i)^2.
\]

Modelo pode ter AUC alta e calibração ruim.

\section{Validação de modelos}

\subsection{Treino, validação e teste}

Treino estima parâmetros. Validação seleciona hiperparâmetros/modelos. Teste é usado uma vez para estimativa final imparcial.

Reutilizar teste repetidamente para tomar decisões transforma-o em validação e cria overfitting ao teste.

\subsection{k-fold cross-validation}

Dados são divididos em $k$ folds. Em cada rodada, $k-1$ treinam e um valida. Métrica final é agregada.

Stratified k-fold preserva proporção de classes aproximadamente em cada fold.

\section{Data leakage}

Leakage ocorre quando informação não disponível no momento real da decisão entra no treino ou pré-processamento.

Exemplos:
\begin{itemize}
\item usar status final do processo para prever risco no início;
\item calcular média de imputação com todo dataset antes de separar teste;
\item selecionar features usando rótulos do teste;
\item misturar registros do mesmo fornecedor em treino/teste quando identificadores quase revelam alvo.
\end{itemize}

\section{Validação temporal}

Quando modelo prevê futuro, divisão temporal é frequentemente mais realista:
\[
\text{treino passado}\rightarrow\text{validação posterior}\rightarrow\text{teste mais recente}.
\]

Random split pode vazar padrões futuros ou entidades repetidas.

\section{Dados ausentes}

Mecanismos conceituais:
\begin{itemize}
\item MCAR: ausência independente de dados observados/não observados;
\item MAR: ausência explicável por dados observados;
\item MNAR: ausência depende do próprio valor ausente ou variável não observada.
\end{itemize}

Imputar média reduz variância e ignora incerteza. Estratégia deve refletir mecanismo e objetivo.

\section{Outliers}

Outlier pode ser:
\begin{itemize}
\item erro de digitação;
\item evento legítimo raro;
\item mudança de regime;
\item fraude;
\item unidade diferente;
\item problema de integração.
\end{itemize}

Excluir automaticamente pode remover justamente o caso de interesse.

\section{Normalização e padronização}

Min-max:
\[
x'=\frac{x-x_{min}}{x_{max}-x_{min}}.
\]

StandardScaler usa média e desvio padrão.

Algoritmos baseados em distância e gradiente podem ser sensíveis à escala. Árvores de decisão, em geral, são invariantes a transformações monotônicas dos atributos usados para split.

\section{Codificação de categóricas}

One-hot evita impor ordem artificial em categoria nominal.

Ordinal encoding é adequado quando ordem possui significado.

Target encoding pode capturar informação útil, mas deve ser calculado dentro de folds de treino para evitar leakage.

\section{Árvores de decisão}

Árvore divide espaço por regras recursivas.

Critérios comuns de classificação incluem Gini e entropia.

Entropia binária:
\[
H=-p\log_2p-(1-p)\log_2(1-p).
\]

Árvore profunda pode overfitar. Controles incluem profundidade, número mínimo de amostras e pruning.

\section{Random Forest}

Combina muitas árvores treinadas em amostras bootstrap e subconjuntos aleatórios de features.

A redução de correlação entre árvores diminui variância do ensemble.

Random forest tende a ser robusta e pouco sensível à escala, mas modelos grandes perdem interpretabilidade direta.

\section{Gradient Boosting}

Boosting adiciona modelos sequencialmente para corrigir erros/resíduos anteriores.

XGBoost, LightGBM e variantes são implementações populares.

Hiperparâmetros como learning rate, número de árvores, profundidade e regularização controlam trade-off viés-variância.

\section{k-NN}

Classifica usando vizinhos mais próximos segundo distância.

É sensível à escala e ao valor de $k$.

$k$ pequeno pode overfitar; $k$ grande suaviza fronteiras e pode underfitar.

Em alta dimensionalidade, distâncias tendem a perder poder discriminativo — curse of dimensionality.

\section{Support Vector Machines}

SVM busca hiperplano de margem máxima. Kernel permite fronteiras não lineares por produtos internos em espaço transformado.

Parâmetro $C$ controla penalização de erros versus margem. Kernel RBF utiliza $\gamma$ associado ao alcance de influência dos pontos.

Escala das features é importante.

\section{Clustering k-means}

K-means minimiza soma das distâncias quadráticas aos centroides:
\[
\sum_{j=1}^k\sum_{x_i\in C_j}\|x_i-\mu_j\|^2.
\]

É sensível a:
\begin{itemize}
\item escala;
\item outliers;
\item inicialização;
\item escolha de $k$;
\item clusters não esféricos.
\end{itemize}

K-means++ melhora inicialização.

\section{DBSCAN}

DBSCAN define regiões densas por $\varepsilon$ e minPts. Pode encontrar formatos irregulares e marcar ruído.

Dificuldades surgem quando clusters têm densidades muito diferentes.

\section{Silhouette}

Para ponto $i$:
\[
s(i)=\frac{b(i)-a(i)}{\max(a(i),b(i))},
\]
onde $a(i)$ mede coesão no cluster e $b(i)$ distância ao cluster vizinho mais próximo.

Valores próximos de 1 indicam boa separação geométrica; isso não prova relevância de negócio.

\section{PCA}

Principal Component Analysis cria combinações lineares ortogonais que maximizam variância explicada.

É técnica não supervisionada e sensível à escala.

Componentes podem reduzir dimensionalidade, mas tornam interpretação dos atributos menos direta.

PCA não seleciona simplesmente ``as features mais importantes''; cria novas variáveis.

\section{Detecção de anomalias}

\subsection{Z-score}

Útil quando distribuição/escala justificam. Limiar como $|z|>3$ é heurístico.

\subsection{IQR}

Robusto para distribuições assimétricas em comparação com média/desvio em certos contextos.

\subsection{Isolation Forest}

Anomalias tendem a ser isoladas por menos divisões aleatórias.

\subsection{Local Outlier Factor}

Compara densidade local de ponto com vizinhos, detectando anomalias contextuais em regiões de densidade diferente.

\section{Regras de associação}

Para regra $A\Rightarrow B$:
\[
support=P(A\cap B),
\]
\[
confidence=P(B\mid A),
\]
\[
lift=\frac{P(B\mid A)}{P(B)}.
\]

Lift $>1$ indica associação positiva relativa à prevalência de B, não causalidade.

\section{Séries temporais}

Série pode conter:
\begin{itemize}
\item tendência;
\item sazonalidade;
\item ciclos;
\item ruído;
\item mudanças estruturais.
\end{itemize}

\subsection{Média móvel}

Suaviza flutuações:
\[
MA_t=\frac{1}{k}\sum_{i=0}^{k-1}x_{t-i}.
\]

\subsection{Autocorrelação}

Observações próximas no tempo podem ser dependentes. Isso viola independência assumida por métodos que tratam registros como iid.

Forecast deve usar validação temporal/rolling origin.

\section{Visualização de dados}

\subsection{Barras}

Comparação de categorias. Eixo deve, em geral, iniciar em zero para comprimento de barras preservar proporção, salvo justificativa explícita.

\subsection{Linhas}

Evolução temporal e tendência.

\subsection{Histograma}

Distribuição de variável contínua. Forma depende da escolha de bins.

\subsection{Boxplot}

Resume mediana, quartis e potenciais outliers.

\subsection{Scatter plot}

Mostra associação entre duas variáveis, clusters e heterocedasticidade.

\section{Gráficos enganosos}

Riscos:
\begin{itemize}
\item eixo truncado;
\item áreas 2D para valores 1D;
\item escala log sem indicação;
\item categorias ordenadas arbitrariamente;
\item cores que sugerem intensidade sem legenda;
\item dupla escala Y que cria correlação visual;
\item esconder denominador.
\end{itemize}

\section{SQL em análise}

SQL deve preservar granularidade.

\subsection{Window functions}

\begin{lstlisting}[language=SQL]
SELECT orgao,
       contrato_id,
       valor,
       SUM(valor) OVER (PARTITION BY orgao) AS total_orgao,
       DENSE_RANK() OVER (
          PARTITION BY orgao
          ORDER BY valor DESC
       ) AS posicao
FROM contrato;
\end{lstlisting}

Window function mantém linhas; \texttt{GROUP BY} as colapsa.

\subsection{JOIN e dupla contagem}

Contrato 1:N Pagamento e 1:N Fiscal. Juntar as duas tabelas diretamente cria produto entre filhos.

A solução é pré-agregar cada relação na granularidade de contrato ou modelar corretamente a medida.

\section{Python e pandas}

\begin{lstlisting}[language=Python]
import pandas as pd

df = pd.read_csv("contratos.csv")

resumo = (
    df.groupby("orgao", as_index=False)
      .agg(valor_total=("valor", "sum"),
           mediana=("valor", "median"),
           contratos=("contrato_id", "nunique"))
)
\end{lstlisting}

Após groupby, granularidade é órgão. Juntar novamente a uma tabela por contrato replicará métricas de órgão; isso pode ser correto para feature contextual, mas não para somar novamente o total.

\section{Reprodutibilidade}

Análise reproduzível registra:
\begin{itemize}
\item origem dos dados;
\item data de extração;
\item query;
\item código;
\item dependências/ambiente;
\item parâmetros;
\item semente aleatória quando relevante;
\item versão do modelo;
\item métricas e gráficos gerados.
\end{itemize}

Notebook editado manualmente sem versão e base sobrescrita dificultam auditoria.

\section{CRISP-DM}

Fases:
\begin{enumerate}
\item entendimento do negócio;
\item entendimento dos dados;
\item preparação;
\item modelagem;
\item avaliação;
\item implantação.
\end{enumerate}

É iterativo. A avaliação pode revelar que definição do alvo está errada e exigir retorno ao entendimento do negócio.

\section{Data analytics em auditoria}

Análises úteis:
\begin{itemize}
\item duplicidade de pagamentos;
\item fracionamento temporal de despesas;
\item concentração de fornecedores;
\item conflitos cadastrais;
\item outliers de preço;
\item pagamentos em dias/horários atípicos;
\item redes de relacionamentos;
\item sequências numéricas anormais;
\item cruzamento com sanções e impedimentos.
\end{itemize}

Resultado analítico é indício que orienta procedimentos adicionais.

\section{Lei de Benford}

Para certos conjuntos de números naturalmente distribuídos em várias ordens de magnitude, primeiro dígito $d$ pode seguir aproximadamente:
\[
P(d)=\log_{10}\left(1+\frac{1}{d}\right),\quad d=1,\ldots,9.
\]

Dígito 1 aparece cerca de 30,1\%; 9 cerca de 4,6\%.

Benford não é apropriada para:
\begin{itemize}
\item números atribuídos, como CPF e protocolos;
\item valores com limites estreitos;
\item preços tabelados;
\item amostras pequenas;
\item dados construídos por regras fixas.
\end{itemize}

Desvio de Benford não prova fraude; é técnica de screening.

\section{Análise de redes}

Grafos representam entidades como nós e relações como arestas.

Métricas:
\begin{itemize}
\item grau: número de conexões;
\item betweenness: frequência em caminhos mínimos;
\item componentes conectados;
\item comunidades;
\item centralidade.
\end{itemize}

Em fraude, uma empresa com grau alto pode ser fornecedora legítima de grande porte. Contexto e evidência adicional são necessários.

\section{Privacidade e minimização na análise}

Nem toda feature disponível deve ser utilizada. Dados pessoais precisam de finalidade e base jurídica, além de controles de acesso e minimização.

A remoção de atributo sensível não elimina proxies. CEP, escola, renda e padrão de consumo podem reconstruir características protegidas.

\section{Síntese conceitual}

\mapafuncional{Análise de Dados}
{Análise válida = amostra adequada + estatística correta + validação realista + interpretação proporcional à evidência}
{Descrição $\rightarrow$ média, mediana, dispersão e gráficos.\\Inferência $\rightarrow$ amostragem, IC e testes.\\Relações $\rightarrow$ correlação e regressão.\\Predição $\rightarrow$ métricas, threshold e validação.\\Exploração $\rightarrow$ clustering, anomalias, regras e séries.\\Auditoria $\rightarrow$ SQL, Benford, redes e reprodutibilidade.}
{Classe rara: accuracy engana.\\p-valor não mede probabilidade de H0.\\AUC mede ranking, não calibração.\\Teste deve imitar produção.\\Outlier é sinal, não conclusão.\\Correlação não identifica direção causal.}
{Amostra grande $\neq$ amostra não enviesada.\\Significância estatística $\neq$ relevância prática.\\R² alto $\neq$ causalidade.\\Benford $\neq$ detector de fraude universal.\\Random split $\neq$ validação adequada para qualquer problema temporal.}
{Qual é a unidade de análise?\\A amostra representa a população?\\Qual o custo de FP e FN?\\A feature existia no momento da decisão?\\O gráfico preserva escala?\\A análise é reproduzível?}

\begin{resumo}
Análise de Dados exige separar descrição, inferência, predição e causalidade. A dificuldade de prova aparece quando uma métrica correta é usada para responder pergunta errada: média em distribuição assimétrica, acurácia em classe rara, p-valor como tamanho de efeito, AUC como calibração ou correlação como causalidade.
\end{resumo}
